\chapter{Black--Scholes from First Principles}
\label{ch:black_scholes}
\chaptermark{Black--Scholes from first principles}

The no-arbitrage bounds of \cref{ch:options_payoffs} pin a fair
European-call price to the interval
$[\max(\Spx_t - \Kstrike e^{-\rrf(\Tmat-t)},\, 0),\,\Spx_t]$. For
an at-the-money one-year call on a $\$100$ stock at $\rrf = 5\%$
that interval is $[4.877,\, 100]$ --- ninety-five dollars wide,
and no portfolio argument narrows it further. A trader needs a
single number. This chapter shows how Black, Scholes, and Merton
(1973) closed the gap by constructing a \emph{self-financing
replicating portfolio} that tracks the option's payoff in every
state. No-arbitrage then forces the option price to equal the
replicating cost --- a single number that depends on the
underlying's \emph{volatility} but not on its drift.

The pre-1973 literature priced options by guessing an expected
payoff under the physical measure and discounting at a risk-%
adjusted rate. Both inputs were unknowable. Black and Scholes'
insight: neither is needed. Continuous stock--bond rebalancing
replicates the payoff exactly, so the price equals the
\emph{cost of replication}. The drift $\mu$ drops out.

\begin{backgroundbox}[Prerequisites]
From earlier in the book:
\begin{itemize}
  \item \Cref{ch:mean_reversion_ou}: SDEs, Brownian motion
  $W_t$, the It\^o integral and isometry. This chapter uses the
  full It\^o lemma rather than just the linear OU
  integrating-factor argument.
  \item \Cref{ch:options_payoffs}: European call/put payoffs,
  no-arb bounds, put--call parity. Black--Scholes refines the
  lower bound $\max(\Spx - \Kstrike e^{-\rrf\Tmat}, 0)$ into a
  single point.
  \item \Cref{ch:arbitrage_zoo}: ``two portfolios with identical
  payoffs in every state must trade at the same price.''
\end{itemize}
From outside the book:
\begin{itemize}
  \item Stochastic calculus at the It\^o-lemma level: if
  $f \in C^{2,1}$ and $X_t$ is an It\^o process with
  $dX_t = a_t\, dt + b_t\, dW_t$, then
  \[
    df(X_t, t) \;=\; \frac{\partial f}{\partial t}\,dt
                 \;+\; \frac{\partial f}{\partial x}\,dX_t
                 \;+\; \tfrac{1}{2}\,\frac{\partial^2 f}{\partial x^2}\,
                       (dX_t)^2,
  \]
  with the It\^o multiplication table
  $(dt)^2 = dt\,dW_t = 0$ and $(dW_t)^2 = dt$.
  \item Heat equation $u_\tau = \tfrac12 u_{xx}$: Gaussian
  fundamental solution, closed-form on $\R$ via convolution
  against initial data.
  \item Standard normal distribution and CDF $\CDFN$.
\end{itemize}
\end{backgroundbox}

\begin{intuitionbox}
Black--Scholes is a hedging argument, not a distributional one.
It does not answer ``what is the expected discounted payoff of
the option?'' --- that question has no model-free answer. It
answers ``how much does it cost to replicate the payoff with
stock and a bond?'' That question has a unique answer once you
fix the volatility, and the answer does not depend on the drift.
\end{intuitionbox}

\section{The Black--Scholes setup}
\label{sec:ch15_setup}

Black--Scholes makes six assumptions. They are restrictive and
often counter-factual, but the formula is robust to moderate
violations of each. List them once and carry them through the
derivation.

\begin{enumerate}
  \item[\textbf{(A1)}] \textbf{Geometric Brownian motion (GBM)}.
  \begin{equation}
    d\Spx_t = \mu\,\Spx_t\,dt + \volat\,\Spx_t\,dW_t,
    \qquad \Spx_0 > 0,
    \label{eq:gbm}
  \end{equation}
  with constant drift $\mu \in \R$, constant volatility
  $\volat > 0$, and $W_t$ standard Brownian under $\Prob$.
  Equivalently $\ln\Spx_t$ is BM with drift, so $\Spx_t$ is
  log-normal. We model \emph{log}-returns (not prices) as
  Gaussian because prices cannot go negative while log-returns
  can, and because a sum of many small independent log-returns
  is approximately normal by the CLT. \emph{Why it matters:}
  gives the closed form $(d\Spx_t)^2 = \volat^2\Spx_t^2\,dt$,
  the only stochastic input It\^o needs.
  \item[\textbf{(A2)}] \textbf{Constant risk-free rate $\rrf$,}
  continuously compounded. \emph{Why:} the bond leg grows
  deterministically as $e^{\rrf t}$ with $dB_t = \rrf B_t\,dt$
  --- no $dW$ contribution.
  \item[\textbf{(A3)}] \textbf{Frictionless trading.} Continuous
  rebalancing, no transaction costs, no spread, no borrow
  cost, no funding, no impact. \emph{Why:} self-financing
  reduces to $d\Pi_t = \Delta_t\,d\Spx_t + \rrf B_t\,dt$ with
  no cash leakage.
  \item[\textbf{(A4)}] \textbf{No dividends} during $[0,\Tmat]$.
  \emph{Why:} long-stock P\&L is only $d\Spx_t$; dividends
  would add a $\Delta_t\,(\text{div})\,dt$ term and modify the
  PDE.
  \item[\textbf{(A5)}] \textbf{No arbitrage:} two portfolios
  with identical payoffs in every state must trade at the same
  price. \emph{Why:} turns ``replicates'' into ``costs the
  same.''
  \item[\textbf{(A6)}] \textbf{European exercise.} Payoff
  $g(\Spx_T)$ only at $\Tmat$; no early exercise. \emph{Why:}
  the boundary condition is imposed at $\Tmat$ alone, not as a
  free boundary at every $t < \Tmat$ (which American options
  require).
\end{enumerate}

\begin{keyfactbox}[Black--Scholes assumptions]
GBM underlying with constant $\volat$ and drift $\mu$; constant
risk-free rate $\rrf$; frictionless continuous trading; no
dividends; no arbitrage; European exercise. Every common
extension (dividends, term-structure of rates, local volatility,
jumps, transaction costs) relaxes exactly one of the six.
\end{keyfactbox}

\begin{pitfallbox}[None of these assumptions are literally true]
Stocks do not follow GBM: returns have fat tails, volatility is
stochastic and time-varying, jumps occur on news, dividends are
paid, rates move, spreads are positive, borrow costs vary, and
you cannot rebalance continuously. Yet Black--Scholes is the
universal language of equity-options markets because traders
read it as a \emph{coordinate system}, not a physical model. Its
job is to translate observed option prices into implied
volatility (\cref{ch:vol_surface}); the resulting surface is the
market's correction to every wrong assumption in (A1)--(A6) at
once.
\end{pitfallbox}

\section{A toy pricing example, by hand}
\label{sec:ch15_binomial}

Before any SDE, price a call by hand on a one-step binomial
tree. This makes the \emph{replication} idea concrete in a
setting with no calculus; every move in the continuous-time
derivation later is a continuous analogue of a move made here in
arithmetic.

\subsection{Setup}
\label{subsec:ch15_binomial_setup}

A non-dividend stock trades at $\Spx_0 = 100$ and over one
period takes either $\Spx_u = 110$ or $\Spx_d = 90$. The
risk-free rate is $\rrf = 2\%$ continuously compounded
($e^{0.02} \approx 1.02020$). Price a European call with
$\Kstrike = 100$. Payoff at expiry:
\begin{align}
  \Call_u \;&=\; \max(\Spx_u - \Kstrike,\, 0)
            \;=\; \max(10,\, 0) \;=\; 10, \nonumber \\
  \Call_d \;&=\; \max(\Spx_d - \Kstrike,\, 0)
            \;=\; \max(-10,\, 0) \;=\; 0.
  \label{eq:bin_payoffs}
\end{align}
Two states, two payoffs. Construct a stock--bond portfolio whose
payoff equals $\Call_u$ in the up-state and $\Call_d$ in the
down-state; no-arbitrage forces it to cost the same as the call.

\subsection{Construction of the replicating portfolio}
\label{subsec:ch15_binomial_construct}

Hold $\Delta$ shares and borrow $B$ dollars (bond leg $-B$
today, debt $Be^{\rrf}$ at expiry). End-of-period value:
\begin{equation}
  \Pi_T \;=\; \Delta\,\Spx_T \;-\; B\,e^{\rrf},
  \label{eq:bin_portfolio}
\end{equation}
with $\Spx_T \in \{\Spx_u, \Spx_d\}$. Requiring $\Pi_T = \Call_T$
in every state gives two equations in two unknowns.

\subsection{The up-state equation}
\label{subsec:ch15_binomial_up}

Substitute $\Spx_T = \Spx_u = 110$:
\begin{equation}
  \underbrace{110}_{=\,\Spx_u}\cdot\,\Delta
  \;\;-\;\;
  \underbrace{B\,e^{0.02}}_{\text{debt repayment}}
  \;\;=\;\;
  \underbrace{10}_{=\,\Call_u}.
  \label{eq:bin_up_eqn}
\end{equation}
One equation, two unknowns.

\subsection{The down-state equation}
\label{subsec:ch15_binomial_down}

Substitute $\Spx_T = \Spx_d = 90$:
\begin{equation}
  \underbrace{90}_{=\,\Spx_d}\cdot\,\Delta
  \;\;-\;\;
  \underbrace{B\,e^{0.02}}_{\text{debt repayment}}
  \;\;=\;\;
  \underbrace{0}_{=\,\Call_d}.
  \label{eq:bin_down_eqn}
\end{equation}
The system is now square.

\subsection{Solve and compute the cost}
\label{subsec:ch15_binomial_solve}

Subtract \eqref{eq:bin_down_eqn} from \eqref{eq:bin_up_eqn}: the
$Be^{0.02}$ terms cancel and we get
\begin{equation}
  20\,\Delta \;=\; 10
  \quad\Longrightarrow\quad
  \boxed{\,\Delta \;=\; 0.5\,}.
  \label{eq:bin_delta}
\end{equation}
Substitute $\Delta = 0.5$ into \eqref{eq:bin_down_eqn}:
\begin{equation}
  90 \cdot 0.5 \;-\; B\,e^{0.02} \;=\; 0
  \quad\Longrightarrow\quad
  B\,e^{0.02} \;=\; 45
  \quad\Longrightarrow\quad
  B \;=\; \frac{45}{e^{0.02}} \;\approx\; 44.10894.
  \label{eq:bin_B}
\end{equation}
The portfolio borrows $\$44.10894$ today. Setup cost at $t = 0$
is the half share minus the borrowed cash:
\begin{align}
  \text{cost}_0
  \;&=\; \Delta\,\Spx_0 \;-\; B
  \;=\; 0.5 \cdot 100 \;-\; 44.10894 \nonumber \\
  \;&=\; 50 \;-\; 44.10894
  \;\approx\; 5.89106. \label{eq:bin_cost}
\end{align}
The portfolio pays exactly the call payoff in both states, so
no-arbitrage gives
\begin{equation}
  \boxed{\;\Call_0 \;=\; 5.89106.\;}
  \label{eq:bin_C0}
\end{equation}

\begin{intuitionbox}
The option costs whatever it takes to replicate its payoff with
stock and a bond. No probabilities appear; we never said how
likely the up-state is. No-arbitrage turns ``does the same'' into
``costs the same.'' Drift, beliefs, expected returns --- none
enter.
\end{intuitionbox}

\subsection{Risk-neutral reinterpretation}
\label{subsec:ch15_binomial_riskneutral}

A second route reaches the same answer. Define the
\emph{risk-neutral probability} $p$ of an up-move as the unique
number making the stock's discounted expected payoff equal its
current price:
\begin{equation}
  \Spx_0 \;=\; e^{-\rrf}\!\left[\,p\,\Spx_u \;+\; (1-p)\,\Spx_d\,\right].
  \label{eq:bin_rn_def}
\end{equation}
Solve for $p$. With $u \defeq \Spx_u/\Spx_0 = 1.1$ and
$d \defeq \Spx_d/\Spx_0 = 0.9$,
\begin{align}
  e^{\rrf} \;&=\; p\,u \;+\; (1-p)\,d \nonumber \\
  p \;&=\; \frac{e^{\rrf} - d}{u - d}
       \;=\; \frac{1.02020 - 0.9}{1.1 - 0.9}
       \;=\; \frac{0.12020}{0.2}
       \;\approx\; 0.6010.
  \label{eq:bin_p}
\end{align}
$p$ is not the real up-probability --- it is whatever makes the
risk-free rate equal the discounted expected stock return. Now
the call's discounted RN expected payoff:
\begin{align}
  \Call_0^{\text{RN}}
  \;&=\; e^{-\rrf}\!\left[\,p\,\Call_u \;+\; (1-p)\,\Call_d\,\right]
  \nonumber \\
  \;&=\; e^{-0.02}\!\left[\,0.6010 \cdot 10 \;+\; 0.3990 \cdot 0\,\right]
  \nonumber \\
  \;&=\; e^{-0.02} \cdot 6.0101
  \;\approx\; 5.89106. \label{eq:bin_C0_rn}
\end{align}
The same number. Replication and risk-neutral expectation are
two views of the same object; \S\ref{sec:ch15_martingale} gives
the continuous-time version.

\begin{keyfactbox}[Discrete-time replication]
For a one-step binomial tree with up-move $\Spx_u$, down-move
$\Spx_d$ and call payoffs $\Call_u, \Call_d$:
\begin{align*}
  \Delta \;&=\; \frac{\Call_u - \Call_d}{\Spx_u - \Spx_d}
  &\text{(replicating share count)}, \\[0.25em]
  p \;&=\; \frac{e^{\rrf} - d}{u - d}
  &\text{(risk-neutral up-probability)}, \\[0.25em]
  \Call_0 \;&=\; e^{-\rrf}\!\left[\,p\,\Call_u + (1-p)\,\Call_d\,\right]
  &\text{(price)}.
\end{align*}
The replication and risk-neutral formulas give the same number
by construction. Drift never appears in either.
\end{keyfactbox}

\subsection{Summary table for the worked example}
\label{subsec:ch15_binomial_summary}

The five steps on one page:

\begin{center}
\begin{tabular}{@{}lll@{}}
\toprule
\textbf{Quantity} & \textbf{Symbol} & \textbf{Value} \\
\midrule
Spot &  $\Spx_0$           & $100$ \\
Up-state stock & $\Spx_u$  & $110$ \\
Down-state stock & $\Spx_d$ & $90$ \\
Strike & $\Kstrike$        & $100$ \\
Risk-free rate (per period) & $\rrf$ & $0.02$ \\
Up-state call payoff & $\Call_u$ & $10$ \\
Down-state call payoff & $\Call_d$ & $0$ \\
\midrule
Replicating share count & $\Delta$ & $(10-0)/(110-90) = 0.5$ \\
Borrowing (debt at $T$) & $B\,e^{\rrf}$ & $0.5 \cdot 90 - 0 = 45$ \\
Borrowing (cash today) & $B$ & $45 / e^{0.02} \approx 44.10894$ \\
\midrule
Replication cost & $\Delta\,\Spx_0 - B$ & $50 - 44.10894 \approx 5.89106$ \\
Risk-neutral up-prob & $p$ & $(e^{0.02} - 0.9)/0.2 \approx 0.6010$ \\
Risk-neutral price & $e^{-\rrf}[p\,\Call_u + (1-p)\,\Call_d]$ & $\approx 5.89106$ \\
\midrule
\textbf{Call price} & $\boxed{\Call_0}$ & $\boxed{\approx 5.89106}$ \\
\bottomrule
\end{tabular}
\end{center}

Both routes agree to all kept digits. Replication cost
\emph{is} the call price by no-arbitrage; the RN expectation is
the same arithmetic re-bracketed.

\subsection{From one step to many: the CRR limit}
\label{subsec:ch15_crr_limit}

The one-step tree is too crude for a real stock. Slice the year
into $N$ steps of length $\Delta t = \Tmat/N$; each step the
stock moves up or down by a small factor. The Cox--Ross--%
Rubinstein (CRR) parameterisation \citep{coxrossrubinstein1979}
chooses
\begin{equation}
  u \;=\; e^{\volat\sqrt{\Delta t}},
  \qquad
  d \;=\; \frac{1}{u} \;=\; e^{-\volat\sqrt{\Delta t}},
  \qquad
  p \;=\; \frac{e^{\rrf\,\Delta t} - d}{u - d},
  \label{eq:crr_params}
\end{equation}
so the per-step log-return has variance $\volat^2\,\Delta t$,
matching GBM \eqref{eq:gbm}. Price by \emph{backward induction}:
write payoffs at terminal nodes $\Spx_0\,u^{N-i}\,d^i$ and roll
back with $\Call_t = e^{-\rrf\,\Delta t}[p\,\Call^u_{t+\Delta t}
+ (1-p)\,\Call^d_{t+\Delta t}]$ at every interior node. Each
node's arbitrage-freeness comes from the same delta-hedge as
\S\ref{subsec:ch15_binomial_construct}.

As $N \to \infty$, the CRR price converges to Black--Scholes
\eqref{eq:bs_call} by the CLT: a sum of $N$ log-binomial
increments is asymptotically Gaussian, so $\ln\Spx_T$ converges
to the log-normal distribution of GBM. Convergence at the
running parameters $\Spx = \Kstrike = 100$, $\rrf = 5\%$,
$\volat = 20\%$, $\Tmat = 1$:

\begin{center}
\begin{tabular}{@{}rrr@{}}
\toprule
$N$ steps & CRR call price & Error vs BS \\
\midrule
$1$    & $12.1623$ & $+1.7117$ \\
$2$    & $9.5405$  & $-0.9101$ \\
$4$    & $9.9705$  & $-0.4801$ \\
$10$   & $10.2534$ & $-0.1972$ \\
$50$   & $10.4107$ & $-0.0399$ \\
$100$  & $10.4306$ & $-0.0200$ \\
$500$  & $10.4466$ & $-0.0040$ \\
$2000$ & $10.4496$ & $-0.0010$ \\
\midrule
$\infty$ (Black--Scholes) & $\mathbf{10.4506}$ & --- \\
\bottomrule
\end{tabular}
\end{center}

Error decays as $O(1/N)$ with even/odd oscillation around ATM.
At $N = 500$ the tree is within four cents of Black--Scholes.
\emph{Black--Scholes is the limit of the discrete
replicating-portfolio price as the rebalance interval shrinks
to zero.}

\begin{intuitionbox}
The Black--Scholes price is the binomial replication price in
the limit of infinitely many infinitesimal rebalances. Every
continuous-time argument below is the binomial argument sent to
$N \to \infty$. If the continuous-time algebra feels mysterious,
drop back to the one-step tree.
\end{intuitionbox}

\section{From binomial to continuous: the PDE derivation}
\label{sec:ch15_pde}

We now repeat the binomial argument in continuous time: form a
stock--bond portfolio, choose the share count to match the call's
dynamics, equate the cost. The new ingredients are It\^o's lemma
and the GBM dynamics \eqref{eq:gbm}; the output is a PDE for
$\Call(\Spx, t)$. We follow the MIT 18.S096 lectures
\citep{mit18S096} and \citeHull{} Ch~13--15.

\subsection{Step 1 --- Set up the hedged portfolio}
\label{subsec:ch15_pde_step1}

Hold $\Delta_t$ shares and a bank-account position $B_t$ at
every $t$. The portfolio's value is
\begin{equation}
  \Pi_t \;=\; \Delta_t\,\Spx_t \;+\; B_t.
  \label{eq:portfolio_value}
\end{equation}
Choose $(\Delta_t, B_t)$ so that $\Pi_t = \Call_t$ for every
$t \in [0, \Tmat]$; by no-arbitrage, $\Call_0 = \Pi_0$.

\subsection{Step 2 --- Self-financing condition}
\label{subsec:ch15_pde_step2}

The portfolio is \emph{self-financing}: cash from sold shares
goes into the bank, cash to buy shares comes out; no external
flows. Instantaneous P\&L then has two sources --- the stock
capital gain $\Delta_t\, d\Spx_t$ and the bond interest
$\rrf\, B_t\, dt$:
\begin{equation}
  d\Pi_t \;=\; \Delta_t\,d\Spx_t \;+\; \rrf\,B_t\,dt.
  \label{eq:self_financing}
\end{equation}
Note what is \emph{absent}: no $\Spx_t\,d\Delta_t$, no $dB_t$
beyond $\rrf B_t\,dt$, no $d\Delta_t\,d\Spx_t$ --- the formal
statement that rebalancing adds no cash.

\subsection{Step 3 --- Apply It\^o's lemma to the call}
\label{subsec:ch15_pde_step3}

We do not yet know $\Call(\Spx, t)$, but assume it is $C^{2,1}$
(verified ex post from the formula). Apply It\^o's lemma:
\begin{equation}
  d\Call_t
  \;=\; \frac{\partial \Call}{\partial t}\,dt
     \;+\; \frac{\partial \Call}{\partial \Spx}\,d\Spx_t
     \;+\; \tfrac{1}{2}\,\frac{\partial^2 \Call}{\partial \Spx^2}\,
           (d\Spx_t)^2.
  \label{eq:ito_C}
\end{equation}
The three terms are time decay, slope in $\Spx$, and the It\^o
curvature correction. The third term is what distinguishes
stochastic from ordinary calculus: smooth paths give
$(d\Spx_t)^2 = O((dt)^2)$ and can be dropped, but Brownian
roughness makes $(d\Spx_t)^2 = O(dt)$ and it must be kept.

The GBM dynamics \eqref{eq:gbm} pin down the quadratic variation:
\begin{align}
  (d\Spx_t)^2
  \;&=\; (\mu\,\Spx_t\,dt + \volat\,\Spx_t\,dW_t)^2
  \nonumber \\
  \;&=\; \mu^2\,\Spx_t^2\,(dt)^2
       \;+\; 2\,\mu\,\volat\,\Spx_t^2\,dt\,dW_t
       \;+\; \volat^2\,\Spx_t^2\,(dW_t)^2
  \nonumber \\
  \;&=\; \underbrace{0}_{(dt)^2 = 0}
       \;+\; \underbrace{0}_{dt\,dW_t = 0}
       \;+\; \volat^2\,\Spx_t^2\,
             \underbrace{(dW_t)^2}_{=\,dt}
  \nonumber \\
  \;&=\; \volat^2\,\Spx_t^2\,dt,
  \label{eq:dS_squared}
\end{align}
using $(dt)^2 = dt\,dW_t = 0$ and $(dW_t)^2 = dt$. Substitute
into \eqref{eq:ito_C}:
\begin{equation}
  d\Call_t
  \;=\;
  \underbrace{\left(\,\frac{\partial \Call}{\partial t}
              \;+\; \tfrac{1}{2}\,\volat^2\,\Spx_t^2\,
                    \frac{\partial^2 \Call}{\partial \Spx^2}\,\right)}_{
              \text{deterministic } dt\text{-coefficient}}
  \,dt
  \;+\;
  \underbrace{\frac{\partial \Call}{\partial \Spx}}_{
              \text{coefficient of }d\Spx_t}
  \,d\Spx_t.
  \label{eq:dC_clean}
\end{equation}
Equation \eqref{eq:dC_clean} is It\^o applied to the call under
GBM --- the only stochastic-calculus input we need
(\citeHull{} Ch~14).

\subsection{Step 4 --- Match the hedged portfolio to the call}
\label{subsec:ch15_pde_step4}

We need $d\Pi_t = d\Call_t$, comparing \eqref{eq:dC_clean} and
\eqref{eq:self_financing}. Two It\^o processes are equal iff
their $dt$- and $dW_t$-coefficients agree separately.

\paragraph{Cancel the stochastic kicks.}
The $dW_t$-coefficient on the call side comes from the
$d\Spx_t = \mu \Spx_t\, dt + \volat \Spx_t\, dW_t$ term in
\eqref{eq:dC_clean}: it is $(\partial\Call/\partial\Spx)\,
\volat\,\Spx_t$. The $dW_t$-coefficient on the portfolio side
comes from the $\Delta_t\, d\Spx_t$ term in
\eqref{eq:self_financing}: it is $\Delta_t\,\volat\,\Spx_t$.
Equating them gives
\begin{equation}
  \boxed{\;\Delta_t \;=\; \frac{\partial \Call}{\partial \Spx}.\;}
  \label{eq:delta_choice}
\end{equation}
\emph{This is the hedge.} Choose the share count to equal the
call's slope in $\Spx$; the Brownian noise in the portfolio
cancels the Brownian noise in the call instant by instant, so
$\Pi_t - \Call_t$ has no $dW_t$ term and is locally riskless.

\begin{intuitionbox}
The stochastic term vanished. Choosing $\Delta_t =
\partial\Call/\partial\Spx$ cancels the portfolio's next
Brownian kick against the call's. What's left is deterministic,
which is why the price ends up depending only on $\volat$
(size of the kicks to hedge) and not on $\mu$ (which drops out
of the $dt$ equation, shown next).
\end{intuitionbox}

\paragraph{Equate the deterministic drifts.}
With $\Delta_t = \partial\Call/\partial\Spx$, the
$dt$-coefficients of \eqref{eq:dC_clean} and
\eqref{eq:self_financing} must also agree. The $dt$-coefficient of
\eqref{eq:dC_clean} is
\begin{equation}
  \frac{\partial \Call}{\partial t}
  \;+\; \tfrac{1}{2}\,\volat^2\,\Spx_t^2\,
        \frac{\partial^2 \Call}{\partial \Spx^2}
  \;+\; \mu\,\Spx_t\,\frac{\partial \Call}{\partial \Spx},
  \label{eq:dt_coeff_C}
\end{equation}
where the last term came from extracting the drift component of
$d\Spx_t$. The $dt$-coefficient of \eqref{eq:self_financing} is
\begin{equation}
  \Delta_t\,\mu\,\Spx_t \;+\; \rrf\,B_t
  \;=\; \mu\,\Spx_t\,\frac{\partial \Call}{\partial \Spx}
       \;+\; \rrf\,B_t,
  \label{eq:dt_coeff_Pi}
\end{equation}
where again the first term is the drift component of
$\Delta_t\,d\Spx_t$ once we substitute $\Delta_t =
\partial\Call/\partial\Spx$. Setting \eqref{eq:dt_coeff_C} equal
to \eqref{eq:dt_coeff_Pi}, the $\mu\,\Spx_t\,\partial\Call/
\partial\Spx$ term cancels from both sides:
\begin{equation}
  \frac{\partial \Call}{\partial t}
  \;+\; \tfrac{1}{2}\,\volat^2\,\Spx_t^2\,
        \frac{\partial^2 \Call}{\partial \Spx^2}
  \;=\; \rrf\,B_t.
  \label{eq:drift_match}
\end{equation}
Now use $\Pi_t = \Call_t$ together with $\Pi_t = \Delta_t\,\Spx_t
+ B_t$ to solve for $B_t$:
\begin{equation}
  B_t \;=\; \Call_t \;-\; \Delta_t\,\Spx_t
       \;=\; \Call_t \;-\; \Spx_t\,\frac{\partial \Call}{\partial \Spx}.
  \label{eq:B_solve}
\end{equation}
Substitute \eqref{eq:B_solve} into \eqref{eq:drift_match}:
\begin{equation}
  \frac{\partial \Call}{\partial t}
  \;+\; \tfrac{1}{2}\,\volat^2\,\Spx_t^2\,
        \frac{\partial^2 \Call}{\partial \Spx^2}
  \;=\; \rrf\,\Call_t
     \;-\; \rrf\,\Spx_t\,\frac{\partial \Call}{\partial \Spx}.
  \label{eq:pde_intermediate}
\end{equation}
Move everything to the left and tidy:
\begin{equation}
  \boxed{\;
    \frac{\partial \Call}{\partial t}
    \;+\; \tfrac{1}{2}\,\volat^2\,\Spx_t^2\,
          \frac{\partial^2 \Call}{\partial \Spx^2}
    \;+\; \rrf\,\Spx_t\,\frac{\partial \Call}{\partial \Spx}
    \;-\; \rrf\,\Call_t
    \;=\; 0.
  \;}
  \label{eq:bs_pde}
\end{equation}

That is the \emph{Black--Scholes PDE}: a linear parabolic PDE in
$(\Spx_t, t)$. The same PDE prices \emph{any} European derivative
whose payoff at $\Tmat$ depends only on $\Spx_T$; the contract
enters only through the boundary condition. For a call,
\begin{equation}
  \Call(\Spx,\,\Tmat) \;=\; (\Spx - \Kstrike)^+,
  \label{eq:bs_bc_call}
\end{equation}
and for a put,
\begin{equation}
  \Put(\Spx,\,\Tmat) \;=\; (\Kstrike - \Spx)^+.
  \label{eq:bs_bc_put}
\end{equation}

\begin{keyfactbox}[The Black--Scholes PDE]
A European derivative on a non-dividend-paying GBM underlying
with constant volatility $\volat$ and constant risk-free rate
$\rrf$ has price $V(\Spx, t)$ satisfying
\[
  \frac{\partial V}{\partial t}
  \;+\; \tfrac{1}{2}\,\volat^2\,\Spx^2\,
        \frac{\partial^2 V}{\partial \Spx^2}
  \;+\; \rrf\,\Spx\,\frac{\partial V}{\partial \Spx}
  \;-\; \rrf\,V
  \;=\; 0,
  \quad
  V(\Spx, \Tmat) = g(\Spx_T),
\]
where $g$ is the contractual payoff. $\mu$ does not appear. The
same PDE prices calls, puts, digitals, and any path-independent
European payoff --- only $g$ changes.
\end{keyfactbox}

\begin{intuitionbox}
$\mu$ disappeared between \eqref{eq:dt_coeff_C} and
\eqref{eq:drift_match}: the same delta-hedge that killed the
$dW_t$ exposure also killed the drift exposure that came along
with it. What survives is the PDE \eqref{eq:bs_pde}, in which
only $\volat$ appears. \emph{A hedged option position does not
care which way the stock drifts; it cares how violently the
stock wiggles}, because the wiggles are what the hedger chases.
Two analysts with different views of $\mu$ agree on the price as
long as they agree on $\volat$.
\end{intuitionbox}

The next section solves \eqref{eq:bs_pde}.

\section{Solving the PDE: change of variables to the heat equation}
\label{sec:ch15_solving}

The Black--Scholes PDE \eqref{eq:bs_pde} is a disguised heat
equation
\begin{equation}
  \frac{\partial u}{\partial \tau}
  \;=\; \tfrac{1}{2}\,\frac{\partial^2 u}{\partial x^2},
  \label{eq:heat}
\end{equation}
solved since Fourier (1822). Three changes of variable remove,
in turn, the multiplicative diffusion, the discount, and the
drift. We work through enough algebra to see the disguise
dissolve, then state the integral form (\citeHull{} Ch~13).

\paragraph{Change of variable 1: log-spot.}
GBM is multiplicative in $\Spx$ but the heat equation is
additive in $x$. Take logs:
\begin{equation*}
  x \;\defeq\; \ln \Spx,
  \qquad
  V(\Spx, t) \;\eqdef\; \widetilde V(x, t).
\end{equation*}
By the chain rule,
\begin{align*}
  \frac{\partial V}{\partial \Spx}
  \;&=\; \frac{1}{\Spx}\,\frac{\partial \widetilde V}{\partial x},
  \\[0.25em]
  \frac{\partial^2 V}{\partial \Spx^2}
  \;&=\; -\,\frac{1}{\Spx^2}\,\frac{\partial \widetilde V}{\partial x}
       \;+\; \frac{1}{\Spx^2}\,\frac{\partial^2 \widetilde V}{\partial x^2}
  \;=\; \frac{1}{\Spx^2}
        \left(\frac{\partial^2 \widetilde V}{\partial x^2}
              \;-\; \frac{\partial \widetilde V}{\partial x}\right).
\end{align*}
Substitute into \eqref{eq:bs_pde}; the $\Spx$ factors cancel
against the chain-rule $1/\Spx^2$ and $1/\Spx$:
\begin{equation}
  \frac{\partial \widetilde V}{\partial t}
  \;+\; \tfrac{1}{2}\,\volat^2
        \left(\,\frac{\partial^2 \widetilde V}{\partial x^2}
              \;-\; \frac{\partial \widetilde V}{\partial x}\right)
  \;+\; \rrf\,\frac{\partial \widetilde V}{\partial x}
  \;-\; \rrf\,\widetilde V
  \;=\; 0.
  \label{eq:bspde_logS}
\end{equation}
The diffusion coefficient is now constant $\tfrac{1}{2}\volat^2$
instead of spot-dependent $\tfrac{1}{2}\volat^2\Spx^2$. Group the
first-order $\partial \widetilde V/\partial x$ terms:
\begin{equation}
  \frac{\partial \widetilde V}{\partial t}
  \;+\; \tfrac{1}{2}\,\volat^2\,\frac{\partial^2 \widetilde V}{\partial x^2}
  \;+\; \bigl(\rrf - \tfrac{1}{2}\volat^2\bigr)\,
        \frac{\partial \widetilde V}{\partial x}
  \;-\; \rrf\,\widetilde V
  \;=\; 0.
  \label{eq:bspde_logS_grouped}
\end{equation}

\paragraph{Change of variable 2: reverse time and remove the discount.}
\eqref{eq:bspde_logS_grouped} runs forward in $t$ with terminal
condition at $\Tmat$; \eqref{eq:heat} is an initial-value problem
in $\tau$. Set
\begin{equation*}
  \tau \;\defeq\; \Tmat - t,
\end{equation*}
so $t = \Tmat$ maps to $\tau = 0$ and
$\partial/\partial t = -\partial/\partial \tau$. Equation
\eqref{eq:bspde_logS_grouped} becomes
\begin{equation}
  -\,\frac{\partial \widetilde V}{\partial \tau}
  \;+\; \tfrac{1}{2}\,\volat^2\,\frac{\partial^2 \widetilde V}{\partial x^2}
  \;+\; \bigl(\rrf - \tfrac{1}{2}\volat^2\bigr)\,
        \frac{\partial \widetilde V}{\partial x}
  \;-\; \rrf\,\widetilde V
  \;=\; 0,
  \label{eq:bspde_taufirst}
\end{equation}
and we can rearrange:
\begin{equation}
  \frac{\partial \widetilde V}{\partial \tau}
  \;=\; \tfrac{1}{2}\,\volat^2\,\frac{\partial^2 \widetilde V}{\partial x^2}
  \;+\; \bigl(\rrf - \tfrac{1}{2}\volat^2\bigr)\,
        \frac{\partial \widetilde V}{\partial x}
  \;-\; \rrf\,\widetilde V.
  \label{eq:bspde_tau}
\end{equation}
\eqref{eq:bspde_tau} is constant-coefficient linear parabolic on
$\R$. It is almost the heat equation but contaminated by a
first-order drift and a zeroth-order discount. Remove both by:
\begin{equation*}
  \widetilde V(x, t) \;\eqdef\; e^{-\rrf\,\tau}\,
                                \widehat V\!\left(x + (\rrf - \tfrac{1}{2}\volat^2)\,\tau,\;\tau\right),
\end{equation*}
i.e.\ pull out $e^{-\rrf\,\tau}$ and shift the spatial variable
along the drift line. The $\rrf\widetilde V$ term cancels
(differentiating $e^{-\rrf\tau}$ yields $-\rrf$), and the
$(\rrf - \tfrac{1}{2}\volat^2)\,\partial\widetilde V/\partial x$
term cancels (the shift removes it). What remains is the heat
equation
\begin{equation}
  \frac{\partial \widehat V}{\partial \tau}
  \;=\; \tfrac{1}{2}\,\volat^2\,\frac{\partial^2 \widehat V}{\partial y^2},
  \qquad y \defeq x + (\rrf - \tfrac{1}{2}\volat^2)\,\tau,
  \label{eq:heat_form}
\end{equation}
on $y \in \R$. Rescaling $y \mapsto y/\volat$ gives the
$\tfrac{1}{2}$-coefficient form \eqref{eq:heat}; the coefficient
is a unit convention, not a methodological distinction.

\paragraph{Solution by Gaussian convolution.}
The heat equation on the real line with initial data $f(y)$ has
the explicit Gaussian-convolution solution
\begin{equation}
  \widehat V(y, \tau)
  \;=\; \int_{-\infty}^{\infty}
        \frac{1}{\volat\sqrt{2\pi\tau}}\,
        e^{-(y - y')^2/(2\volat^2 \tau)}\,
        f(y')\,dy'.
  \label{eq:gaussian_convolution}
\end{equation}
Fourier's 1822 formula in Black--Scholes coordinates. The call
boundary $(\Spx - \Kstrike)^+$ rewritten in $y$, with the four
substitutions undone, becomes an integral of
$(\Spx_T - \Kstrike)^+$ against the lognormal density of
$\Spx_T$ under $\Qm$ --- the integral
\S\ref{sec:ch15_martingale} sets up directly. Evaluation:
split at $\Kstrike$, complete the square, recognise two
Gaussian tails (\citeHull{} Ch~13;
\S\ref{subsec:ch15_lognormal_eval} below).

\begin{intuitionbox}
The Black--Scholes PDE is a heat equation in disguise: the
underlying diffuses multiplicatively (cured by $x = \ln\Spx$),
the price is discounted (cured by $e^{-\rrf\tau}$), and the
risk-neutral drift shifts the distribution (cured by
translating $y$). After the three substitutions we are
computing the diffusion of heat from the payoff kink
$(\Spx - \Kstrike)^+$ over time-to-expiry $\Tmat - t$; the
smoothing of the kink is the option's time value.
\end{intuitionbox}

\section{The Black--Scholes formula}
\label{sec:ch15_formula}

Solving \eqref{eq:bs_pde} with the call boundary
\eqref{eq:bs_bc_call} gives:

\begin{keyfactbox}[Black--Scholes formula]
The price at $t$ of a European call on a non-dividend GBM stock
(strike $\Kstrike$, maturity $\Tmat$, constant $\volat$,
constant $\rrf$) is
\begin{equation}
  \boxed{\;
    \Call(\Spx_t,\,t)
    \;=\;
    \Spx_t\,\CDFN(d_1)
    \;-\;
    \Kstrike\,e^{-\rrf(\Tmat - t)}\,\CDFN(d_2),
  \;}
  \label{eq:bs_call}
\end{equation}
where $\CDFN$ is the standard normal CDF and
\begin{align}
  d_1 \;&=\;
  \frac{\ln(\Spx_t/\Kstrike)
        \;+\; \bigl(\rrf + \tfrac{1}{2}\volat^2\bigr)\,(\Tmat - t)}
       {\volat\,\sqrt{\Tmat - t}},
  \label{eq:bs_d1} \\[0.4em]
  d_2 \;&=\;
  d_1 \;-\; \volat\,\sqrt{\Tmat - t}.
  \label{eq:bs_d2}
\end{align}
The corresponding put price, derivable from \eqref{eq:bs_call}
by put--call parity (\cref{ch:options_payoffs}), is
\begin{equation}
  \Put(\Spx_t,\,t)
  \;=\;
  \Kstrike\,e^{-\rrf(\Tmat - t)}\,\CDFN(-d_2)
  \;-\;
  \Spx_t\,\CDFN(-d_1).
  \label{eq:bs_put}
\end{equation}
\end{keyfactbox}

Sanity check: differentiating \eqref{eq:bs_call} twice in
$\Spx$ and once in $t$, substituting into \eqref{eq:bs_pde},
reduces it to $0 = 0$ (we verified in sympy; \citeHull{} Ch~13
reproduces the algebra).

\subsection{Interpreting each term}
\label{subsec:ch15_terms}

The two terms have a clean economic reading. The first
$\Spx_t\,\CDFN(d_1)$ is the risk-neutral expected present value
of receiving the stock at expiry, conditional on expiring ITM.
The second $\Kstrike\,e^{-\rrf(\Tmat - t)}\,\CDFN(d_2)$ is the
expected PV of paying the strike, conditional on exercising.
The call is ``what you receive minus what you pay,'' both
conditional on exercise.

The two CDF arguments encode the two relevant probabilities:
\begin{itemize}
  \item $\CDFN(d_2) = \Qm[\Spx_T \geq \Kstrike]$ is the
  risk-neutral ITM probability --- the probability you have to
  pay the strike, hence its appearance in the second term.
  \item $\CDFN(d_1)$ is the call's \emph{Delta},
  $\partial\Call/\partial\Spx$, the same $\Delta$ as in
  \eqref{eq:delta_choice}. \Cref{ch:greeks_hedging} derives it
  by direct differentiation.
\end{itemize}

\begin{intuitionbox}
$\CDFN(d_2)$ is the ITM probability; $\CDFN(d_1)$ is the Delta.
Memorise both. Desks read $\CDFN(d_2)$ as ``probability of
printing ITM at expiry (RN)'' and $\CDFN(d_1)$ as ``share count
to delta-hedge.'' Those are the two facts that turn the formula
into desk language.
\end{intuitionbox}

\subsection{Worked numerical example}
\label{subsec:ch15_numerical}

We use the same running parameters as the no-arb bound at the
start of the chapter:
\begin{equation*}
  \Spx \;=\; 100, \qquad \Kstrike \;=\; 100,
  \qquad \rrf \;=\; 5\%, \qquad \volat \;=\; 20\%,
  \qquad \Tmat - t \;=\; 1 \text{ year.}
\end{equation*}

\paragraph{Step 1 --- Compute $d_1$ and $d_2$.}
With $\ln(\Spx/\Kstrike) = \ln 1 = 0$ and $\Tmat - t = 1$,
\begin{align}
  d_1 \;&=\; \frac{0 + (0.05 + \tfrac{1}{2}\cdot 0.20^2)\cdot 1}
                  {0.20 \cdot \sqrt{1}}
       \;=\; \frac{0.05 + 0.02}{0.20}
       \;=\; \frac{0.07}{0.20}
       \;=\; 0.35,
  \label{eq:numex_d1} \\[0.25em]
  d_2 \;&=\; d_1 \;-\; \volat\,\sqrt{\Tmat - t}
       \;=\; 0.35 \;-\; 0.20
       \;=\; 0.15.
  \label{eq:numex_d2}
\end{align}

\paragraph{Step 2 --- Standard normal CDF.}
$\CDFN(0.35) \approx 0.6368$, $\CDFN(0.15) \approx 0.5596$.

\paragraph{Step 3 --- Discount factor.}
$e^{-\rrf(\Tmat - t)} = e^{-0.05} \approx 0.95123$.

\paragraph{Step 4 --- Plug into the call formula.}
\begin{align}
  \Call \;&=\; \Spx\,\CDFN(d_1)
            \;-\; \Kstrike\,e^{-\rrf(\Tmat - t)}\,\CDFN(d_2)
  \nonumber \\
  \;&=\; 100 \cdot 0.6368
       \;-\; 100 \cdot 0.95123 \cdot 0.5596 \nonumber \\
  \;&=\; 63.683 \;-\; 53.232 \nonumber \\
  \;&\approx\; 10.4506.
  \label{eq:numex_call}
\end{align}
The ATM one-year call is worth $\approx \$10.45$ per share.

\begin{example}[Sanity-check against the no-arb interval]
\Cref{ch:options_payoffs} gave $\Call \in [4.877,\, 100]$ for
these parameters. Black--Scholes' $10.4506$ sits inside. The
lower bound is the $\volat \to 0$ limit of \eqref{eq:bs_call};
turning vol on lifts the price by the time value
$10.45 - 4.877$. The upper bound is the $\volat \to \infty$
limit. Black--Scholes places a single point inside the
model-free interval.
\end{example}

\paragraph{Step 5 --- Parity check.}
Put--call parity gives $\Call - \Put = \Spx -
\Kstrike\,e^{-\rrf(\Tmat - t)} = 100 - 95.123 = 4.877$, so
$\Put \approx 5.574$. Computing the put from \eqref{eq:bs_put}
directly gives the same number.

\subsection{The replicating portfolio at $t = 0$}
\label{subsec:ch15_replicating_continuous}

At $t = 0$, $\Delta_0 = \CDFN(d_1) = \CDFN(0.35) \approx 0.6368$
shares. The bond position from \eqref{eq:B_solve}:
\begin{equation*}
  B_0 \;=\; \Call_0 \;-\; \Delta_0\,\Spx_0
       \;=\; 10.4506 \;-\; 0.6368 \cdot 100
       \;\approx\; -53.23,
\end{equation*}
so the portfolio is long $0.6368$ shares ($\$63.68$) and short
$\$53.23$ in cash. Net cost $63.68 - 53.23 \approx 10.45$
matches the call price. These are the continuous analogues of
the binomial $\Delta = 0.5$, $B = 44.10894$, and cost $5.89106$.

The hedge is dynamic: as $\Spx_t$ moves and time decays both
$\Delta_t = \CDFN(d_1(\Spx_t, t))$ and $B_t$ change.
\Cref{ch:greeks_hedging} computes the rebalance rate (Gamma) and
the discrete-rebalance P\&L. For now: the replicating portfolio
at every $(\Spx_t, t)$ is a deterministic function of two
variables --- no hidden state.

\subsection{Comparative statics: how the price moves with each input}
\label{subsec:ch15_comparative}

The price depends on five inputs: $\Spx, \Kstrike, \rrf, \volat,
\Tmat - t$. Holding four fixed at the base case ($\Spx =
\Kstrike = 100$, $\rrf = 5\%$, $\volat = 20\%$, $\Tmat - t = 1$)
and perturbing one:

\begin{center}
\begin{tabular}{@{}rr@{\quad}rr@{\quad}rr@{}}
\toprule
\multicolumn{2}{c}{\textbf{Vary spot}}
& \multicolumn{2}{c}{\textbf{Vary vol}}
& \multicolumn{2}{c}{\textbf{Vary $\Tmat - t$}} \\
$\Spx$ & $\Call$ & $\volat$ & $\Call$ & $\Tmat - t$ & $\Call$ \\
\midrule
$80$  & $1.86$  & $10\%$ & $6.80$  & $0.10$ & $2.77$  \\
$90$  & $5.09$  & $20\%$ & $10.45$ & $0.25$ & $4.62$  \\
$100$ & $10.45$ & $30\%$ & $14.23$ & $0.50$ & $6.89$  \\
$110$ & $17.66$ & $40\%$ & $18.02$ & $1.00$ & $10.45$ \\
$120$ & $26.17$ & $50\%$ & $21.79$ & $2.00$ & $16.13$ \\
\bottomrule
\end{tabular}
\end{center}
\noindent\emph{In each column, only the named variable changes;
the other four are held at the base case $\Spx = \Kstrike =
100$, $\rrf = 5\%$, $\volat = 20\%$, $\Tmat - t = 1$.}

The qualitative facts every options trader knows:
\begin{itemize}
  \item \textbf{Monotone increasing in spot.}
  $\partial\Call/\partial\Spx = \CDFN(d_1) = \Delta \in (0, 1)$,
  with $\Delta \to 1$ deep ITM, $\to 0$ deep OTM. Spot-column
  slope near $\Spx = 100$ is $(17.66-5.09)/20 \approx 0.63$,
  matching $\CDFN(0.35) = 0.6368$.
  \item \textbf{Monotone increasing in vol.} Higher vol widens
  the terminal distribution; the convex call payoff turns the
  upside into more expected value than the downside takes away.
  $\partial\Call/\partial\volat$ is \emph{Vega}.
  \item \textbf{Monotone increasing in time-to-expiry.} More
  time means more opportunity to drift ITM. $\partial\Call/
  \partial t$ is \emph{Theta} (negative as $t \to \Tmat$).
  \item \textbf{Increasing in the rate.} Higher $\rrf$ pulls
  the discounted strike down, raising ATM/ITM prices.
  Sensitivity is \emph{Rho} --- small short-dated equity, large
  long-dated rates.
\end{itemize}
All four are derived in \cref{ch:greeks_hedging}; see
\citeHull{} Ch~19 for textbook treatment.

\subsection{A reference Python implementation}
\label{subsec:ch15_code}

A 12-line Python function. Every options-trading interview asks
the candidate to write this; it should be reflex.

\begin{lstlisting}[caption={Black--Scholes European call price as a pure function. Returns a single float; no side effects, no plotting, no persistence. Exactly the function a Strat would type into a notebook to sanity-check a vendor pricer.}]
from math import log, sqrt, exp
from scipy.stats import norm

def black_scholes_call(S, K, r, sigma, T):
    """Price of a European call on a non-dividend-paying GBM
    underlying with constant rate and constant vol.
        S     spot
        K     strike
        r     continuously compounded risk-free rate
        sigma volatility (annualised)
        T     time to maturity in years
    Returns the call price in the same units as S.
    """
    if T <= 0.0:
        return max(S - K, 0.0)
    d1 = (log(S / K) + (r + 0.5 * sigma ** 2) * T) / (sigma * sqrt(T))
    d2 = d1 - sigma * sqrt(T)
    return S * norm.cdf(d1) - K * exp(-r * T) * norm.cdf(d2)
\end{lstlisting}

Calling \texttt{black\_scholes\_call(100, 100, 0.05, 0.20, 1.0)}
returns $\approx 10.4506$, matching \eqref{eq:numex_call}. The
$T \leq 0$ branch handles the expiry boundary; production code
also branches on $\volat = 0$ to avoid division by zero.

\section{A second derivation: the martingale route}
\label{sec:ch15_martingale}

The PDE derivation is the 1973 route and makes the hedge
explicit. A second route due to Harrison--Pliska (1981) reaches
the same formula via a measure change. We sketch it because both
routes matter in practice and because the martingale route is
the foundation of most numerical methods in
\cref{ch:numerical_methods}.

\subsection{The risk-neutral measure}
\label{subsec:ch15_riskneutral_measure}

Under physical $\Prob$, GBM \eqref{eq:gbm} reads
$d\Spx_t = \mu\,\Spx_t\,dt + \volat\,\Spx_t\,dW_t^{\Prob}$, so
the discounted stock $\widetilde\Spx_t \defeq e^{-\rrf t}\Spx_t$
has drift $(\mu - \rrf)\widetilde\Spx_t\,dt$ and is generally
not a martingale. Girsanov constructs an equivalent measure
$\Qm$ under which $\widetilde\Spx_t$ \emph{is} a martingale.
(Equivalent measures agree on which events have probability
zero but disagree on the weight each event carries.) Under
$\Qm$, $\Spx_t$ has dynamics
\begin{equation}
  d\Spx_t \;=\; \rrf\,\Spx_t\,dt
              \;+\; \volat\,\Spx_t\,dW_t^{\Qm},
  \label{eq:gbm_Q}
\end{equation}
where $W_t^{\Qm}$ is $\Qm$-Brownian. The drift under $\Qm$ is
$\rrf$, not $\mu$; the volatility is unchanged (Girsanov rotates
drift only \citeHull{}). In plain English, $\Qm$ is the unique
reweighting of outcomes under which every tradable asset earns
the risk-free rate on average --- the investor behind $\Qm$ is
indifferent to risk, so replacing $\mu$ with $\rrf$ is exactly
what eliminates the risk premium that hedging already neutralised.
$\Qm$ is the \emph{risk-neutral measure} --- the continuous
analogue of the binomial $p$ from \eqref{eq:bin_p} (\citeHull{}
Ch~28).

\subsection{Pricing as a discounted expectation}
\label{subsec:ch15_pricing_expectation}

Once $\widetilde\Spx_t$ is a $\Qm$-martingale, every discounted
self-financing portfolio is too --- in particular
$e^{-\rrf t}\Call_t$. The martingale property
$\E^{\Qm}[e^{-\rrf\Tmat}\Call_T \mid \mathcal F_t] =
e^{-\rrf t}\Call_t$ gives
\begin{equation}
  \boxed{\;
    \Call_t \;=\; e^{-\rrf(\Tmat - t)}\,
                  \E^{\Qm}\!\left[\,(\Spx_T - \Kstrike)^+
                                  \;\big|\; \mathcal F_t\,\right].
  \;}
  \label{eq:risk_neutral_pricing}
\end{equation}
The price is the discounted risk-neutral expected payoff. No
$\mu$, no preferences, no ``true'' ITM probability appears; the
inputs are $\Spx_t, \rrf, \volat, \Kstrike, \Tmat$.

\subsection{Evaluating the lognormal expectation}
\label{subsec:ch15_lognormal_eval}

Under $\Qm$ the dynamics \eqref{eq:gbm_Q} integrate to
\begin{equation}
  \Spx_T \;=\; \Spx_t\,\exp\!\left(\,
    \bigl(\rrf - \tfrac{1}{2}\volat^2\bigr)(\Tmat - t)
    \;+\; \volat\,\sqrt{\Tmat - t}\,Z\,\right),
  \qquad Z \sim \Normal(0, 1),
  \label{eq:ST_lognormal}
\end{equation}
(the It\^o-lemma solution applied to $\ln\Spx_t$). So $\Spx_T$
is log-normal with log-mean $\ln\Spx_t + (\rrf - \tfrac12\volat^2)
(\Tmat - t)$ and log-variance $\volat^2(\Tmat - t)$. The
expectation is an integral of the payoff against the lognormal
density:
\begin{equation*}
  \E^{\Qm}\!\left[(\Spx_T - \Kstrike)^+\right]
  \;=\; \int_{-\infty}^{\infty}
        \left(\Spx_t\,e^{(\rrf - \volat^2/2)(\Tmat-t)
                          + \volat\sqrt{\Tmat-t}\,z} - \Kstrike\right)^+
        \,\frac{e^{-z^2/2}}{\sqrt{2\pi}}\,dz.
\end{equation*}
The integrand is non-zero for $z > -d_2$. Splitting at $-d_2$
and evaluating (first piece: complete the square; second:
Gaussian tail) gives
\begin{equation*}
  \E^{\Qm}[(\Spx_T - \Kstrike)^+]
  \;=\; \Spx_t\,e^{\rrf(\Tmat - t)}\,\CDFN(d_1)
       \;-\; \Kstrike\,\CDFN(d_2).
\end{equation*}
Multiplying by $e^{-\rrf(\Tmat - t)}$ recovers \eqref{eq:bs_call}
exactly. \citeHull{} Ch~14 walks the integral line by line;
\citep{mit18S096} gives a parallel treatment. We confirmed the
equivalence numerically: both give $10.4506$ at the running
parameters to ten significant figures.

\begin{intuitionbox}
The PDE route says ``solve a deterministic PDE driven by the
payoff''; the martingale route says ``compute a discounted
expectation under a measure that prices the discounted stock
as a martingale.'' The \emph{Feynman--Kac theorem} formalises
the equivalence: any parabolic PDE
$\partial_t V + \mathcal L V = 0$ with terminal condition
$V(\Spx, \Tmat) = g(\Spx)$ equals
$\E[g(\Spx_T) \mid \Spx_t]$ under the diffusion with generator
$\mathcal L$. Black--Scholes is Feynman--Kac for risk-neutral
GBM. \Cref{ch:numerical_methods} uses both: finite differences
for the PDE, Monte Carlo for the expectation.
\end{intuitionbox}

\section{What each assumption buys and when it breaks}
\label{sec:ch15_assumptions_break}

The six assumptions are not equally restrictive: some cost a
few basis points, others are catastrophic on a news event. The
implied-vol surface (\cref{ch:vol_surface}) is the market's
simultaneous correction to all six. Below: which fix lives
where.

\paragraph{(A1) Constant volatility breaks first.}
Real volatility clusters, spikes on news, and varies with
moneyness (the \emph{smile/skew}). The market quotes a different
implied vol per strike--maturity. Stochastic-vol models (Heston,
SABR) make $\volat$ its own SDE; local-vol (Dupire) makes
$\volat$ a deterministic function of $(\Spx, t)$ fitting today's
surface. \Cref{ch:equity_derivatives} covers both.

\paragraph{(A1) Lognormality breaks too.}
Real returns have fat tails (kurtosis $\gg 3$), negative skew,
and discontinuous jumps. Black--Scholes is a continuous
diffusion and cannot model jumps. Jump-diffusion (Merton 1976,
variance-gamma) and rough-vol models address the gap, but the
desk response is simpler: trade the implied-vol surface, which
absorbs the model error into the strike-dependent vol.

\paragraph{(A2) Constant rates: fine for short-dated equity, fatal for rates derivatives.}
For an SPX option with $\Tmat - t < 1$ year, rate variation
gives sub-cent pricing error vs the spread. For a 30-year cap
or swaption it is absurd: the contract is \emph{about} the
random rate. \Cref{ch:rates_derivatives} covers the proper
models (Black 1976, Hull--White, BGM, SABR).

\paragraph{(A3) Frictionless trading breaks at the rebalance frequency.}
You cannot rebalance continuously. With minimum interval
$\Delta t$ the hedge has $O(\sqrt{\Delta t})$ discretisation
error and positive replication-P\&L variance. The bid--ask
spread adds a cost linear in turnover. \Cref{ch:greeks_hedging}
computes the per-step discrete-hedge P\&L as
$\tfrac12\Gamma(\Delta\Spx)^2 - \Theta\,\Delta t$; expected zero
(gamma cancels theta on average), realised noisy and
concentrated in the gamma term --- the basis of gamma scalping.

\paragraph{(A4) No dividends --- clean fix.}
For known discrete dividends $D_i$ at times $t_i$, replace
$\Spx_t$ in \eqref{eq:bs_call} with $\Spx_t - PV(D)$ where
$PV(D) = \sum_i D_i\,e^{-\rrf(t_i - t)}$. For a continuous yield
$q$, replace $\Spx_t$ with $\Spx_t\,e^{-q(\Tmat - t)}$ and use
$\rrf - q$ in $d_1$. \citeHull{} Ch~17 covers both.

\paragraph{(A6) European exercise rules out American puts.}
American puts on dividend payers (or on non-payers with positive
rates) can be optimal to exercise early; \eqref{eq:bs_call} does
not price them. \Cref{ch:numerical_methods} treats American
options via finite-difference PDEs with the early-exercise
inequality and via Longstaff--Schwartz regression.

\begin{pitfallbox}[Don't calibrate Black--Scholes to historical vol and trade off it]
Estimating $\volat$ from a year of historical returns, plugging
into \eqref{eq:bs_call}, and trading against market price loses
money systematically:
\begin{itemize}
  \item \textbf{Historical vol is backward-looking; the market
  is forward-looking.} A Tesla call two days before earnings is
  priced off expected vol, not realised. Earnings, mergers,
  FOMC, launches all show up in implied vol first.
  \item \textbf{Different strikes trade at different implied
  vols.} A deep-OTM put has much higher IV than an ATM call on
  the same name, pricing in crash asymmetry. ``The vol'' as a
  scalar is a category error; the unit is the \emph{surface}.
\end{itemize}
The professional response inverts: treat \eqref{eq:bs_call} as a
function from price to \emph{implied volatility}, invert
numerically per strike--maturity, and trade the surface against
your view (\cref{ch:vol_surface}).
\end{pitfallbox}

\begin{prosperitybox}
In Prosperity 3, Round 3 introduced
\texttt{VOLCANIC\_ROCK\_VOUCHERS} --- five European calls on
\texttt{VOLCANIC\_ROCK}, strikes $9500$--$10500$ in steps of
$250$, underlying near $10000$ \citep{frankfurthedgehogs2025}.
Each voucher granted the right to buy one rock at strike at
expiry, with seven days to expiry at the start of Round 3
dropping to two by the final round.

The synthetic options had every real-call feature: positive time
value, monotone-decreasing in strike, delta-hedgeable. Top-10
teams (Frankfurt Hedgehogs, Linear Utility, others) priced with
\eqref{eq:bs_call} and traded short-horizon IV dislocations
(\cref{ch:vol_surface}): buy cheap vol, sell rich vol,
delta-hedge.

Concrete number: at Round 4 start, $\Spx = 10000$. Take the
$\Kstrike = 9750$ voucher with $\Tmat - t \approx 7/365 \approx
0.01918$ and a recent-history $\volat = 20\%$. With $\rrf = 0$
(no risk-free rate in Prosperity),
\begin{align*}
  d_1 \;&=\; \frac{\ln(10000/9750) + (0 + 0.02)\cdot 0.01918}
                  {0.20 \cdot \sqrt{0.01918}}
       \;\approx\; 0.928, \\
  d_2 \;&=\; d_1 - 0.20 \cdot \sqrt{0.01918}
       \;\approx\; 0.900, \\
  \Call \;&=\; 10000 \cdot \CDFN(0.928)
            \;-\; 9750 \cdot \CDFN(0.900)
       \;\approx\; 276.8.
\end{align*}
Intrinsic is $\max(\Spx - \Kstrike, 0) = 250$, so time value is
$\approx \$26.8$ for one week of optionality on a 20\%-vol
underlying. Frankfurt Hedgehogs fitted the IV smile across the
five strikes, computed each market price's residual from the
smile, and traded mean reversion of the residual
\citep{frankfurthedgehogs2025} --- impossible without
evaluating \eqref{eq:bs_call} on every tick.
\end{prosperitybox}

\begin{gsbox}
On a GS equity-derivatives desk, Black--Scholes is the baseline
pricer for every listed European option. The Strat's job is not
to derive it but to \emph{calibrate the IV surface}
(\cref{ch:vol_surface}) so that \eqref{eq:bs_call} at each
calibrated $(\volat, K, T)$ reproduces market quotes. The
pricing system is a Black--Scholes loop over the surface; the
calibration loop sits above, adjusting the surface to fit.

Risk is computed by perturbing the IV surface, not the
``physical'' vol. Sensitivities (vega per strike, vega per
tenor, cross-vega) come from bumping surface nodes and
re-pricing. Traders ask ``what is my vega to the 1y 25-delta
put?'', not ``what is sigma?'' The Strat owns the calibration
and bumping pipeline; the trader owns the position.
\Cref{ch:risk_production} expands.
\end{gsbox}

\begin{historybox}
Black and Scholes published ``Pricing of Options and Corporate
Liabilities'' in May 1973 in the \emph{JPE} after multiple
rejections (``too narrowly applied to finance''). Merton
independently derived the formula in a companion paper, with
the rigorous self-financing and no-arbitrage treatment. The
CBOE had opened one month earlier; within a few years floor
traders carried programmable calculators preloaded with the
formula. Black died in 1995; Scholes and Merton received the
1997 Nobel, with the committee noting Black's contribution.
\end{historybox}

\section{Key facts to memorise}
\label{sec:ch15_keyfacts}

\begin{itemize}
  \item \textbf{Setup}: GBM
  $d\Spx_t = \mu\Spx_t\,dt + \volat\Spx_t\,dW_t$, constant
  $\rrf$, frictionless continuous trading, no dividends, no
  arbitrage, European exercise.
  \item \textbf{Binomial replication}: $\Delta =
  (\Call_u - \Call_d)/(\Spx_u - \Spx_d)$ plus a bond from the
  down-state equation track the call. Risk-neutral
  $p = (e^{\rrf} - d)/(u - d)$ gives the same price as a
  discounted expectation. Drift never appears.
  \item \textbf{Black--Scholes PDE}:
  \[
    \frac{\partial \Call}{\partial t}
    + \tfrac{1}{2}\volat^2\Spx^2
          \frac{\partial^2 \Call}{\partial \Spx^2}
    + \rrf\Spx\frac{\partial \Call}{\partial \Spx}
    - \rrf\Call = 0,
  \]
  with terminal $\Call(\Spx, \Tmat) = (\Spx - \Kstrike)^+$.
  \item \textbf{Call formula}:
  \[
    \Call = \Spx\,\CDFN(d_1) - \Kstrike\,e^{-\rrf(\Tmat - t)}\,
    \CDFN(d_2),
  \]
  with $d_1 = [\ln(\Spx/\Kstrike) + (\rrf + \tfrac12\volat^2)
  (\Tmat - t)]/(\volat\sqrt{\Tmat - t})$ and
  $d_2 = d_1 - \volat\sqrt{\Tmat - t}$.
  \item \textbf{Put}:
  $\Put = \Kstrike\,e^{-\rrf(\Tmat - t)}\CDFN(-d_2)
   - \Spx\,\CDFN(-d_1)$ (or via parity).
  \item \textbf{Interpretations}: $\CDFN(d_2)$ = risk-neutral
  ITM probability; $\CDFN(d_1)$ = Delta = replicating share
  count.
  \item \textbf{Numerical anchor}: $\Spx = \Kstrike = 100$,
  $\rrf = 5\%$, $\volat = 20\%$, $\Tmat = 1$ gives
  $\Call \approx 10.4506$, $\Put \approx 5.5735$, parity
  $\Call - \Put \approx 4.877$.
  \item \textbf{$\mu$ does not appear.} Two analysts with
  different drift views agree on price as long as they agree on
  $\volat$ --- the central economic content.
  \item \textbf{Two routes, same answer}: PDE (replicate, kill
  the $dW$, heat equation) and martingale (change to $\Qm$,
  discounted expectation, lognormal integral). Feynman--Kac
  makes the equivalence an identity.
  \item \textbf{No assumption is literally true.} Vol is
  stochastic, returns have fat tails and jumps, hedging is
  discrete and costly, dividends are paid. The market's
  simultaneous correction is the IV surface of
  \cref{ch:vol_surface}.
\end{itemize}
